\subsection{PART A – Review Questions}

\subsubsection{Define the following terms in your own words}

\noindent\textbf{Organization}: 

An organization is a structured group of people working together to achieve common goals through a formal hierarchy, roles, and processes.

\noindent\textbf{Business}: 

A business is a profit-driven organization engaged in commercial, industrial, or professional activities, aiming to deliver goods or services to customers.

\noindent\textbf{Professionalism}: 

Professionalism refers to the conduct, attitude, and standards expected from individuals in a professional setting, including ethics, accountability, competence, and respect.

\noindent\textbf{Organizational value}: 

Organizational value is the worth an initiative (like IT) brings to a company by enhancing efficiency, effectiveness, customer satisfaction, or achieving strategic objectives.

\noindent\textbf{IT investment}: 

An IT investment involves the allocation of resources (time, money, personnel) to information technology systems, tools, or processes to support and enhance business operations.

\subsubsection{List 3 key differences between a business and an organization.}

\indent

\begin{table}[h]
\centering
\begin{tabular}{|c|l|l|}
\hline
\textbf{Aspect}              & \multicolumn{1}{c|}{\textbf{Business}}                                            & \multicolumn{1}{c|}{\textbf{Organisation}}                                               \\ \hline
\textbf{Primary Goal}        & Profit-driven                                                                     & Profit or purpose-driven                                                                 \\ \hline
\textbf{Legal Entity}        & \begin{tabular}[c]{@{}l@{}}Typically \\ a commercial entity\end{tabular}          & \begin{tabular}[c]{@{}l@{}}Can be non-commercial \\ (NGOs, schools, etc.)\end{tabular}   \\ \hline
\textbf{Scope of Activities} & \begin{tabular}[c]{@{}l@{}}Focused on product/services \\ for market\end{tabular} & \begin{tabular}[c]{@{}l@{}}Includes social, educational, \\ religious goals\end{tabular} \\ \hline
\end{tabular}
\end{table}

\subsection{PART B – Group Discussion}

\subsubsection{How do different organizational structures (functional, matrix, flat, hierarchical) impact?}

\noindent\textbf{Communication across departments}: 

\begin{itemize}
    \item \textbf{Functional}: Clear within departments, but \textbf{siloed} between them.
    \item \textbf{Matrix}: \textbf{Complex communication} due to dual reporting.
    \item \textbf{Flat}: \textbf{Open and direct}, promoting transparency.
    \item \textbf{Hierarchical}: \textbf{Top-down} flow, can be \textbf{slow and rigid}.
\end{itemize}

\noindent\textbf{IT alignment with business goals}: 

\begin{itemize}
    \item \textbf{Functional}: Alignment may be \textbf{limited} to department-level priorities.
    \item \textbf{Matrix}: Better alignment as IT may be involved in \textbf{multiple projects}.
    \item \textbf{Flat}: \textbf{Greater collaboration}, easier alignment.
    \item \textbf{Hierarchical}: Requires \textbf{strong top-level direction} to align IT.
\end{itemize}

\noindent\textbf{Decision making in IT investments}: 

\begin{itemize}
    \item \textbf{Functional}: Often \textbf{centralized}, may overlook broader needs.
    \item \textbf{Matrix}: \textbf{Shared decision-making}, balancing multiple interests.
    \item \textbf{Flat}: \textbf{Faster, collective} decisions.
    \item \textbf{Hierarchical}: \textbf{Slower}, but structured and \textbf{top-down}.
\end{itemize}

\subsubsection{Why is it important for IT professionals to understand the business model and operating model of an organization?}

\begin{itemize}
    \item Understanding the \textbf{business model} helps IT professionals see\textbf{ how value is created and delivered}, allowing them to tailor IT solutions accordingly.
    \item Knowing the \textbf{operating model} ensures IT aligns with \textbf{core processes}, \textbf{decision-making structures}, and strategic goals, leading to more \textbf{effective IT investments} and improved \textbf{business outcomes}.
\end{itemize}

\subsection{PART C – Case Study Discussion}

SmartCare Health Ltd. is a growing private healthcare provider operating across three Australian states. The company has recently adopted digital patient management systems and is considering further investment in AI-based diagnostics and cloud-based infrastructure. They currently operate under a divisional structure, with each state operating semi-independently. However, IT systems are centrally managed. The CEO is keen on aligning IT more strategically with the business, but the CFO is hesitant, stating that IT is a cost centre with unclear returns.

\bigskip

\noindent\textbf{What type of organisational structure is SmartCare using?}

A \textbf{divisional structure}, where each \textbf{state operates semi-independently}, but IT is \textbf{centrally managed}.

\noindent\textbf{What potential issues might arise with centralised IT and divisional operations?}

\begin{itemize}
    \item \textbf{Misalignment} between local needs and central IT solutions.
    \item \textbf{Delayed responses} to local issues due to \textbf{central control}.
    \item \textbf{Resistance} from divisions feeling \textbf{disconnected} from IT decisions
\end{itemize}


\noindent\textbf{How might aligning IT strategy with business goals create organisational value for SmartCare?}

\begin{itemize}
    \item Enhances \textbf{patient outcomes} through data integration and AI insights.
    \item Improves \textbf{efficiency and scalability} across divisions.
    \item Enables \textbf{strategic decision-making} with data-driven tools.
    \item Supports \textbf{innovation} and \textbf{competitive advantage} in healthcare
\end{itemize}

\noindent\textbf{Suggest one-way SmartCare could measure the value of its IT investment in the AI diagnostics system}

Track diagnostic accuracy improvements, reduced diagnosis time, and cost savings using KPIs like turnaround time, patient throughput, and error rates pre- and post-implementation.
